\section{From Physical to Serviceable Capacity}
\subsection{Capacity and Access Demand}
Let $x_{jp}$ denote bytes of class $j$ resident in pool $p$, and $C_p$
its usable capacity. For a specified observation period $\Delta$, let
$a_j$ be the fraction of resident bytes accessed and $A_{jp}$ the
applicable transfer amplification. A simplified demand model is
\begin{align}
  \sum_j x_{jp} &\le C_p, \\
  D_p = \frac{\sum_j x_{jp}a_jA_{jp}}{\Delta}
       &\le B_{p,\mathrm{useful}}.
  \label{eq:demand}
\end{align}
All terms must use consistent byte definitions. Read requests, responses,
writes, and on-wire overhead require separate resource accounting.
Constraints also apply to each die's HBM and each shared path; multiple
pools cannot each receive the full budget of the same bottleneck link.
Equation~\ref{eq:demand} is a necessary average-rate condition, not a
sufficient deadline or tail-latency guarantee.

\subsection{Distinct Concurrency Lifetimes}
An internal MPU slot is occupied from memory admission until data service
finishes. An endpoint transaction can remain outstanding until the
response reaches its specified acceptance or consumption point. Network
credits have their own allocation and return events. These lifetimes
must not be conflated.

For homogeneous transactions of $g$ useful bytes, a diagnostic upper
bound is
\begin{equation}
\begin{split}
 B_{p,\mathrm{useful}} \le \min\{&B_{\mathrm{mem}}, B_{\mathrm{path}},
 B_{\mathrm{recv}},\\
 &N_s g/L_s,\; N_e g/R_e,\; P u\},
\end{split}
\label{eq:bound}
\end{equation}
where $N_s$ and $L_s$ describe service slots and their mean occupancy,
$N_e$ and $R_e$ describe endpoint transactions, and $P$ and $u$ describe
packet throughput and useful bytes per packet. Terms must be converted
to useful-byte rates. The bound is a diagnostic approximation rather
than a new performance law; queueing, dependencies, and credit return
require explicit timing models.

\begin{figure}[t]
\centering
\setlength{\fboxsep}{6pt}
\fbox{\begin{minipage}{0.91\columnwidth}
\centering
Per-die placement and remote admission\\[3pt]
$\downarrow$ request path / shared fabric\\[3pt]
MPU admission $\rightarrow$ latency $\rightarrow$ data service\\
\emph{Internal service slot released here}\\[3pt]
$\downarrow$ response queue / shared fabric\\[3pt]
Finite receiver buffer $\rightarrow$ HBM consumption\\
\emph{Proposed endpoint budget released at consumption}
\end{minipage}}
\caption{Proposed control boundary. Internal memory slots, network
credits, and endpoint transaction budgets have separate lifetimes.
This is a mechanism schematic, not a measured execution trace.}
\label{fig:lifetime}
\end{figure}

An illustrative calculation uses the current configuration's 128 service
slots, 256-byte requests, and 200-ns fixed memory latency:
\begin{equation}
  \frac{128\times256\;\mathrm{B}}{200\;\mathrm{ns}}
  = 163.84\;\mathrm{GB/s}.
\end{equation}
This loose upper bound is below the configured 190-GB/s memory bandwidth
even before data service and queueing extend occupancy. It is arithmetic
from assumed parameters, not a hardware measurement. Approximately 149
slots cover only the fixed latency at that rate; they do not guarantee
full throughput. Substituting network round-trip time into this internal
slot calculation would misrepresent the current model.

\subsection{Metric Definition}
For a fixed access class, time budget, and hardware allocation, we define
\emph{serviceable capacity} as the largest resident working set satisfying
the specified completion budget. Mixed classes produce a feasible region,
not a universal scalar capacity. Ultimately, deployment value should be
measured by completed requests meeting both TTFT and TPOT objectives
under a fixed external arrival stream, with rejected work reported.
Representative-stage simulation alone cannot establish those metrics.
